% !TEX root = tarasov-qasp-zh.tex
% 章首扉图与格言（sec-10，场论部分）在英文源 ch-21.tex 尾部（有效代码），归 zh/ch-21-zh.tex；
% 本章文件首尾均无有效 sec 块（ch-22.tex 尾部那组 sec 块为 % 注释的死代码副本，未搬入）。

\chapter[我们来谈谈场论]{我们来谈谈场论\label{ch-22}}

\begin{dialogue}

\tchr 我们来谈谈场（field）——基本的物理概念之一。为了明确起见，我们将只讨论静电场（electrostatic field）。你对场有什么认识？你认为它究竟是什么？

\stda 我得承认，对于场究竟是什么，我的认识十分模糊。场是一种难以捉摸的、看不见的东西，有点像幽灵。可同时又有人说，它弥漫于整个空间。把场定义为一种物质实体（material entity），我并不反对，但这对我毫无意义。说到物质（matter）时，我明白我们谈的是什么；说到场，我就束手无策了。

\stdb 在我看来，场的概念是完全具体的。在任何实物（substance）中，物质可以说是以凝聚的形式存在的；而在场中，物质则可以说是“弥散”在整个空间。我们用肉眼看不见场，这什么也证明不了。借助相当简单的仪器，场可以“看”得非常清楚。场起着物体间相互作用传递者的作用。例如，静电场传递静止电荷之间的相互作用。可以说，每个电荷都在自己周围建立一个场。一个电荷建立的场作用于另一个电荷；反过来，第二个电荷建立的场又作用于第一个电荷。电荷间的库仑（静电）相互作用就是这样实现的。

\stda 可是，难道就不能不用任何“中间人”吗？是什么妨碍我们设想一个电荷直接作用于另一个电荷呢？

\stdb 你的设想会引起严重的异议。假定由于某种原因，其中一个电荷在某一时刻移动了一下（即“挪动”了）。若从“直接相互作用”的设想出发，我们就不得不作出结论：第二个电荷必须在同一瞬间也“挪动”一下。这就意味着，来自第一个电荷的信号在一瞬间就到达第二个电荷。这与相对论（theory of relativity）的基本原理相矛盾。而如果存在相互作用的传递者，即场，信号就通过场从一个电荷传播到另一个电荷。无论传播速度多大，它总是有限的。因此会存在这样一段时间间隔：第一个电荷已经停止“挪动”，而第二个电荷还没有开始挪动。在这段时间里，“挪动”的信号只存在于场之中。
\pnote{学生乙这段论证，正是近代物理放弃“超距作用”（action at a distance）观念的理由。精确地说，电磁相互作用扰动的传播速度就是真空中的光速 $c=299\,792\,458\ \mathrm{m/s}$（精确值；1983 年起米由光速定义）：一个电荷动了，相距 $r$ 的另一个电荷要到 $r/c$ 时间之后才“知道”；瞬时作用意味着无穷大的信号速度，为相对论所不容。场中扰动作波状传播、以光速行进，这正是麦克斯韦电磁波理论的图像（人教版必修第三册第十三章《电磁感应与电磁波初步》）。附带一提：2017 年双中子星并合事件（GW170817）中，引力波与伽马射线暴从约 1.3 亿光年外传到地球仅相差约 1.7 秒，直接证实引力相互作用同样以光速传播——“场传递相互作用”并不限于电磁作用。}

\stda 尽管如此，我还是想听一个关于场的精确定义。

\tchr 我怀着很大兴趣听了你们的对话。我觉得\textbf{学生乙}对现代物理的问题表现出浓厚的兴趣，读过各类物理科普书籍，从而形成了一种不妨称为“主动性思维”的东西。对他来说，场是一个十分现实的、“用得上”的概念；他关于场是相互作用传递者的议论完全正确。\textbf{学生甲}呢，看来只停留在对教科书的形式化阅读上，结果他的思维在相当程度上效率不高。我这样说当然无意冒犯他或任何人，只是想指出：许多考生在类似场合会感到十分无助。说来奇怪，相当多的学生几乎从不读任何科普读物。不过，还是让我们回到问题的实质上来吧。（对\textbf{学生甲}）你要求给场下一个精确定义；没有这样的定义，场的概念就从你身边溜走了。可是你刚才说，你明白什么是物质。那么，你当真知道物质的精确定义吗？

\stda 物质这个概念不需要什么定义。物质是可以用手“摸”到的。

\tchr 既然如此，场的概念同样“不需要什么定义”；场也是可以“摸”到的，只是不用手去摸罢了。然而，定义问题实际上要严肃得多。给出一个精确的、逻辑上无懈可击的定义，意味着用一些更“原始”的概念来表达所定义的概念。可是，如果给定的概念恰好就是“原始”概念之一，又该怎么办呢？不信就试着给几何中的直线下个定义。物质和场这两个概念的情形大致相同：它们都是原始的、基础性的概念，很难指望为它们找到明确的、包罗万象的定义。

\stda 尽管这样，我们总还是能找到某个多少说得通的定义吧？

\tchr 当然能。只是必须记住，任何这样的定义都不可能是详尽无遗的。物质可以以不同的形态存在。它可以集中在空间的一个有限区域内，边界或多或少是确定的（也就是说，它是“定域”（localized）的）；反过来，它也可以是“非定域”（delocalized）的。物质的这些形态中的第一种，可以同“实物”意义上的物质概念联系起来，第二种则同“场”的概念联系起来。两种形态除各自的独特性质外，还具有一些共同的物理特征。例如，有实物单位体积的能量，也有场单位体积的能量；可以谈实物单位体积的动量，也可以谈场单位体积的动量。每一种场传递一种确定的相互作用，而恰恰正是通过这种相互作用，我们才能确定场在任意需要之点的特性。举例来说，带电体在周围空间建立静电场。要显示这个场并测出它在空间某点的强度（intensity），必须把另一个带电体带到该点，并测出作用在这个带电体上的力。这里假定第二个带电体足够小，以致它引起的场的畸变可以忽略不计。
\pnote{用“另一个足够小的带电体”探测场、由测得的力来表征场，正是电场强度的定义方法：人教版必修第三册第九章第 3 节《电场 电场强度》把试探电荷（test charge）放入电场中某点，以 $E=F/q$ 定义该点的电场强度，方向为正电荷受力的方向，并同样要求试探电荷足够小，以免改变被探测的场。原书此处只说“强度”，下一章将把它正式定义为矢量。}

物质的性质不可穷尽，认识的过程永无止境。在认识和实际运用我们周围物质的性质的道路上，我们一步一步地逐渐前进。前进途中，我们不时地要“贴标签”，这些标签是通往知识之路上的某种路标。现在我们把某个东西标记为“场”。我们明白，这个“某个东西”实际上是一片原始深渊。关于这个被我们称作“场”的深渊，我们已经知道得很多，因而能够或多或少令人满意地使用这个新引入的概念。我们知道得很多，但远非全部。想给它一个明确的定义，就如同想丈量无底深渊的深度。
\pnote{“场单位体积的能量与动量”并非修辞：电磁场携带着能量和动量，这是“场也是物质的一种形态”最直接的物理证据。光是电磁波，照射到物体上被吸收或反射时会对物体施加压强（光压）；太阳帆正是利用光压在太空中推进，2010 年日本的 IKAROS 探测器首次成功演示。场的能量密度、动量密度中学不展开，大学电磁学中由麦克斯韦理论严格给出。}

\stdb 我认为，场这个概念，同我们在研究物质世界过程中产生的一切其他概念一样，是不可穷尽的。正因为如此，才不可能给场下一个详尽无遗的、明确的定义。

\tchr 我完全同意你的看法。

\stda 您关于实物和场是物质的两种形态——定域的和非定域的——的论述，我完全满意。可您为什么要扯出关于物理概念的不可穷尽性和认识的永恒性这一大篇议论呢？这些话一入耳，清晰感又消失了，一切又变得有点模糊、有点含混。

\tchr 我理解你的心境。你想寻求的是关于场的某种安稳的定义，哪怕它并不绝对精确；你愿意把它认真地背下来，一经要求就能端出来。你不愿意承认，情况根本不是静止的，而恰恰相反是发展变化的。你不应当认为一切变得模糊了、含混了；我要说，一切变得动态了，即处于趋向变化的状态。任何精确的定义，其本身都是某种僵硬的、终局性的东西，而物理概念恰恰应当在其发展状态中去研究。昨天我们所理解的场概念，同今天我们在这个名称下所理解的已有明显差别。比如，现代物理与经典物理不同，不再在场和实物之间划出一条截然的界限。在现代物理中，场和实物是相互转化的：实物可以变为场，场也可以变为实物。不过，现在就详细讨论这个问题，未免走得太远了。
\pnote{教师预告的“实物与场相互转化”在现代物理中有确切的机制：光子是电磁场的量子，正、负电子对湮没时转化为光子（实物变场），能量足够大的光子反过来又能产生正、负电子对（场变实物）；量子场论进一步把每种基本粒子都看作相应场的激发，“粒子”与“场”不再是两类并立的存在。中学阶段在光的粒子性（光子）与质能方程、质量亏损部分会遇到这类转化的具体例子（人教版选择性必修第三册第四章《原子结构和波粒二象性》、第五章《原子核》）。}

\stdb 我们这场物理讨论明显转到了哲学上。

\tchr 这很自然：任何关于物理概念的讨论，都必然要求参加者具备相当发达的辩证思维（dialectical thinking）能力。这种能力还没有养成时，我们就不得不违背本意地做些哲学性质的迂回。正因为如此，我一再建议你们读各种各样的书，越多越好。这样你们就训练了自己的思维器官，使它更灵活、更生动。在这方面，列宁（V. I. Lenin）的《唯物主义和经验批判主义》（\emph{Materialism and Empiriocriticism}）一书能给任何年轻人以可贵的帮助。我建议你们读一读。
\pnote{这部书出版于 1909 年，写作于物理学刚经历剧变的年代：X 射线（1895）、放射性（1896）、电子（1897）相继发现，“物质消灭了”的说法一度流行。列宁在书中回应的正是这类说法：物质是标志客观实在的哲学范畴，它不依赖于意识而存在；物理学的新发现推翻的只是关于物质具体构造的旧图景，而不是物质本身。本章“物质与场是无法精确定义的原始概念、其内涵随科学发展而不断丰富”的讨论，接续的正是这场争论。}

\stda 可那是一本很难的书，是高等院校的学生研读的。

\tchr 我并不坚持要你们研读这本书。它当然不是供轻松消遣的读物。只要试着认真把它通读一遍。视你们的基础而定，这本书对你思维方式的影响会有大有小。无论如何，它是有益的。

最后我想指出一点：学生甲显然害怕模糊、害怕不确定，他要求最大限度的精确。他忘记了凡事都有一个合理的限度，精确也不例外。试设想一个完全精确的世界，我们对它拥有完备的信息。你就把这样一个世界想象出来，然后告诉我：面对它的原始简单、它的再也无法发展，你难道不吃惊吗？把这一切好好想想，先别急着下结论。现在，暂时让我们试着从另一个角度接近这个问题。我提出如下的问题：“场是怎样描述的？”我知道，许多人得到答案之后会说：“这下我们知道场是什么了。”
\pnote{对“拥有完备信息的完全精确的世界”的否定，在量子力学中有确切的形式：不确定关系表明，粒子的位置与动量不可能同时具有任意精确的值，“对世界拥有完备而精确信息”的图像（拉普拉斯式的决定论）已被现代物理排除。中学阶段不展开，这是大学普通物理中量子力学的基本原理之一。}

\end{dialogue}

\cleardoublepage

% NEWTERMS: field | 场 | 全书首现（章题“场论”与本章正文），首现括注英文；人教版必修第三册第九章起系统展开
% NEWTERMS: electrostatic field | 静电场 | 本章首现括注；ch-23 起系统展开
% NEWTERMS: field theory | 场论 | 章题用语
% NEWTERMS: substance | 实物 | 与场并列为物质的两种（存在）形态；勿与“物态”（state of matter）混淆
% NEWTERMS: material entity | 物质实体 | 学生甲引述的定义用语，照译
% NEWTERMS: transmitter of interactions / go-between | 相互作用的传递者／“中间人” | 原书口语比喻照译
% NEWTERMS: direct interaction | 直接相互作用 | 即通行名“超距作用”（action at a distance），评注给出
% NEWTERMS: localized / delocalized (unlocalized) | 定域／非定域 | 原书加引号的拟语，首现括注；底本另有一处异写 unlocalized，同译不注
% NEWTERMS: intensity (of a field) | （场的）强度 | 本章探测语境首现括注；ch-23 正式定义为电场强度 $E$
% NEWTERMS: theory of relativity | 相对论 | 本章首现括注（泛称，语境为狭义相对论的信号速度上限）
% NEWTERMS: dialectical thinking | 辩证思维 | 照译
% NEWTERMS: test charge | 试探电荷 | 评注层引入（人教版 9.3 术语）；原文作“另一个足够小的带电体”
% NEWTERMS: action at a distance | 超距作用 | 评注层给出通行名
% NEWTERMS: light pressure | 光压 | 评注层引入
% NEWTERMS: solar sail | 太阳帆 | 评注层引入（IKAROS，2010）
% NEWTERMS: photon | 光子 | 评注层引入
% NEWTERMS: annihilation / pair production | 湮没／（电子对）产生 | 评注层引入
% NEWTERMS: quantum field theory | 量子场论 | 评注层引入
% NEWTERMS: uncertainty relation | 不确定关系 | 评注层引入（海森伯）；中学不要求
% NEWTERMS: Materialism and Empiriocriticism | 《唯物主义和经验批判主义》 | 列宁著，1909 年出版；用通行汉译名，正文括注英文原名
